
\documentclass[12pt,a4paper]{report}

% ==== PAQUETES BÁSICOS ====

\usepackage{subcaption}
\usepackage[utf8]{inputenc}
\usepackage{amsmath,amssymb}
\usepackage[T1]{fontenc}
\usepackage[spanish]{babel}
\usepackage{graphicx} % para imágenes
\usepackage{caption}
\usepackage{hyperref} % hipervínculos
\usepackage{geometry} % márgenes
\usepackage{booktabs}
\usepackage{subcaption}
\geometry{left=3cm,right=2.5cm,top=3cm,bottom=2.5cm}

% ==== CÓDIGO FUENTE ====
\usepackage[linesnumbered,ruled,vlined]{algorithm2e}
\SetKwInOut{KwIn}{Entrada}
\SetKwInOut{KwOut}{Salida}

\usepackage{listings}
\usepackage{xcolor}
\usepackage{tikz}
%%%%%Uso de bonificador de javascript
\usepackage{listings}
\usepackage{xcolor}

\lstdefinestyle{javascript}{
  language=Java,
  morekeywords={async,await,import,export,from},
  keywordstyle=\color{blue}\bfseries,
  stringstyle=\color{red},
  commentstyle=\color{gray},
  basicstyle=\ttfamily\small,
  breaklines=true,
  showstringspaces=false
}
%%%%%%%%%%%%%%%%%%%%%%%%%%%%%%%%


\usetikzlibrary{shapes.geometric, arrows,positioning}
\tikzstyle{startstop} = [rectangle, rounded corners, minimum width=2cm, minimum height=0.7cm,text centered, draw=black, fill=red!30]
\tikzstyle{process} = [rectangle, minimum width=2cm, minimum height=0.7cm, text centered, draw=black, fill=blue!20]
\tikzstyle{decision} = [diamond, minimum width=3cm, minimum height=0.7cm, text centered, draw=black, fill=green!30]
\tikzstyle{arrow} = [thick,->,>=stealth]
\lstset{
    language=C,                 % Puedes cambiarlo: Python, Java, etc.
    basicstyle=\ttfamily\small,   % Fuente monoespaciada
    keywordstyle=\color{blue},    % Palabras clave
    stringstyle=\color{red},      % Strings
    commentstyle=\color{gray},    % Comentarios
    backgroundcolor=\color{black!5},
    numbers=left,                 % Números de línea
    numberstyle=\tiny\color{gray},
    stepnumber=1,
    breaklines=true,
    tabsize=4
}



\begin{document}
% ==== PORTADA ====
\begin{titlepage}
    \centering
    
    % Imagen institucional (ejemplo: logo)
    \includegraphics[width=0.3\textwidth]{ug.png}\par\vspace{1cm}
    
    % Nombre de la universidad
    {\Large \textbf{Universidad de Guanajuato}}\\[0.5cm]
    
    % Nombre de la materia
    {\large \textbf{Materia: Aprendizaje de máquina}}\\[1cm]
    {\large \textbf{Carrera: Licenciatura en Ingeniería de Datos e Inteligencia Artificial}}\\[3cm]
    
    % Título de la tarea
    {\Huge \textbf{Tarea 2: Uso de aprendizaje no supervisado para identificación de tareas}}\\[2cm]
    
    % Datos del alumno
    {\large \textbf{Alumno:} José Adrián Rodríguez González}\\[0.3cm]
    {\large \textbf{NUA:} 148661}\\[1.5cm]
    
    % Datos del profesor
    {\large \textbf{Profesor:} Dr. Dora Almanza Luz Ojeda}\\[3cm]
    
    % Fecha
    {\large \today}
    
\vfill
\end{titlepage}

\tableofcontents % Índice automático
\newpage

% ==== INTRODUCCIÓN ====
\chapter{Introducción y objetivos}
En este proyecto se realizó y siguió el esquema completo de toma de conjunto de datos, preprocesamiento de datos, y modelado de IA.
Para esta ocasión se atiende al conjunto de datos de las posiciones de camas. Son dos experimentos. En el segundo experimento se estudian dos casos. Al tomarse y descragarse os conjuntos de datos se preprocesa las imagenes y se calculan las características de las imágenes por medio de la técnica de HoG. Finalmente se implementan diversos modelos de machine learning, abarcando todos los revisados en clase exceptuando modelos de aprendizaje no supervisado.
\chapter{Revisión de los conjuntos de datos}
La motivación del desarrollo de este conjunto de datos, es debido a que la mayoría de los problemas de salud que se agravan a la vejez, suelen deberse a la calidad del sueño. Por eso se hace enfásis en el análisis de una buena calidad de sueño. Y el estudio comienza desde una parte más elemental, como es el como las personas se recuestan. Para ello, se utilizan camas con sensores de presión, que miden cuanta presión hay dentro de una región. De esa manera, se puede dibujar una imagen de la persona y determinar cual es  la posición que tomo al recostarse.
Para el experimento 1, hubieron 13 participantes de los que s eles tomó la captura del sensor en 8 posiciones estandar (la referencia de LA IMAGEN Y LA IAMGEN AQUI).
Para el experimento 1, cada paciente 17 archivos de texto que contienen en cada línea la imagen de $64 \times 32$ falttenizada, dando un fila con 2048 columnas.
Cada archivo contiene tamaños diferentes de tomas o muestras 
Para el experimento 2.Se tomaron más posiciones de cada lado, sin embargo se hizo con dos tipos de camas diferentes. de aire y de esponja.
Cada cama tiene un tamaño de queen size.
(MUestra la imagen que se tomo y su REFERENCIA).Para el caso de las imágenes tomadas den el experimento 2, cada paciente tiene su carpeta de imagenes con las que tiene los dos tipos de camas y en el nombre se adjunta el nombre de la posición tomada, (REFERENCIA DE LA IMAGNE Y SU nombre).

\section {separación de lso conjuntos de datos}
Para ello se debió de implementar un script en python que permita la lectura de las carpetas y la lectura de los archivos.Con el experimento 1, la organizaciṕn de las carpetas es más sencillo, ay que en al raiz de esa carpeta, contiene las carpetas de los pacientes y entor contiene las imágenes.
Para facilitar la clasificación, ya que fueron varias posiciones tomadas por paciente, y en cada posición no tienen la misma cantidad de información, y para no convertir el problema de con demasiadas clases, se decie utilizar solamente 3 clases, posiciones que estén a la izquierda, a la derecha y al centro.
Al separar cada posición acorde a lo qe se ve en (referencia DE LA IMAGEN), se calculan las muestras por clase y hay menos meustras para el lado izquierdo y para el lado derecho. por lo que el proceso siguiente fue hacer una aumentación de datos. Con la aumentación de los datos, cada clase tendría 10314 muestras.
Para el caso del experimento 2, primero se separo por tipo de cama, y el proceso de la toma de muestra fue similiar, sin embargo, la diferencia s que en el arcivo e texto, la matriz de la imagen de presión ya no se encotraba falttenizado, sino que cada columna y renglón corresponde al tamaño original de la imagen.

\section{Aumnetación de los datos}
Se utilizaron pocos filtros y que dieran poca varaición al conjunto de la imagenes originales, ya que la diferencia entre el labelde las imagenes de la derecha y de la izquierda era de 6000 imágenes, por lo ue las iagenes que e ncecesitaban tenían que ser de menor variación posible.
\bibliographystyle{ieeetr}
\bibliography{reference}

\end{document} 
